% =========================================================================================
% glossary.tex -- acronyms and technical terms
%
% This replaces the `glossaries` package, whose \printnoidxglossary sorts entries with TeX
% macros and was slow enough to exhaust Overleaf's free compile budget across the three
% passes of a full build.
%
% Each entry is declared once and used twice: \gls{key} expands it in the text, and the
% same declaration accumulates a row for the printed list. Entries print in declaration
% order, so keep each block alphabetical by hand.
% =========================================================================================

\makeatletter

\newcommand{\acronymrows}{}
\newcommand{\termrows}{}

% \newabbr{key}{SHORT}{long form}
\newcommand{\newabbr}[3]{%
  \expandafter\gdef\csname glsshort@#1\endcsname{#2}%
  \g@addto@macro\acronymrows{\item[#2] #3}%
}

% \newterm{key}{Name}{description}
\newcommand{\newterm}[3]{%
  \expandafter\gdef\csname glsshort@#1\endcsname{#2}%
  \g@addto@macro\termrows{\item[#2] #3}%
}

% \gls prints the short form. An undefined key renders in red rather than silently
% printing nothing, so a typo is visible in the PDF.
\newcommand{\gls}[1]{%
  \ifcsname glsshort@#1\endcsname\csname glsshort@#1\endcsname
  \else\rmq{[undefined term: #1]}\fi}
\newcommand{\glspl}[1]{\gls{#1}s}

% leftmargin/labelwidth must exceed the widest label, or long ones (Grad-CAM, AUROC,
% "Catastrophic forgetting") spill onto their own line and the list looks ragged.
\newcommand{\printacronymlist}{%
  \begin{description}[style=nextline, leftmargin=3.4cm, labelwidth=3.0cm,
                      font=\bfseries, itemsep=0.15em]
    \acronymrows
  \end{description}}

\newcommand{\printtermlist}{%
  \begin{description}[style=nextline, leftmargin=4.6cm, labelwidth=4.2cm,
                      font=\bfseries, itemsep=0.35em]
    \termrows
  \end{description}}

\makeatother

% =============================================================== ACRONYMS (front matter) ==
\newabbr{ai}{AI}{Artificial Intelligence}
\newabbr{auroc}{AUROC}{Area Under the Receiver Operating Characteristic curve}
\newabbr{bo}{BO}{Business Objective}
\newabbr{camus}{CAMUS}{Cardiac Acquisitions for Multi-structure Ultrasound Segmentation}
\newabbr{cnn}{CNN}{Convolutional Neural Network}
\newabbr{crisp}{CRISP-DM}{Cross-Industry Standard Process for Data Mining}
\newabbr{cv}{CV}{Cross-Validation}
\newabbr{dso}{DSO}{Data Science Objective}
\newabbr{dvt}{DVT}{Deep Vein Thrombosis}
\newabbr{ece}{ECE}{Expected Calibration Error}
\newabbr{ed}{ED}{Emergency Department}
\newabbr{ef}{EF}{Ejection Fraction}
\newabbr{fast}{FAST}{Focused Assessment with Sonography for Trauma}
\newabbr{gradcam}{Grad-CAM}{Gradient-weighted Class Activation Mapping}
\newabbr{llm}{LLM}{Large Language Model}
\newabbr{lv}{LV}{Left Ventricle}
\newabbr{mvp}{MVP}{Minimum Viable Product}
\newabbr{pocus}{POCUS}{Point-of-Care Ultrasound}
\newabbr{rag}{RAG}{Retrieval-Augmented Generation}
\newabbr{tta}{TTA}{Test-Time Augmentation}

% ================================================================== TERMS (back matter) ==
\newterm{alines}{A-lines}{Horizontal artefacts repeating at regular intervals below the
  pleural line, produced by normally aerated lung.}

\newterm{balancedacc}{Balanced accuracy}{The mean of the per-class recalls. Under class
  imbalance it is the honest alternative to plain accuracy, which a model can inflate by
  predicting the majority class.}

\newterm{blines}{B-lines}{Vertical bright artefacts extending from the pleural line to the
  bottom of a lung ultrasound image. Their number and distribution indicate fluid in the
  lung interstitium, and they are the principal sonographic sign of pulmonary oedema.}

\newterm{calibration}{Calibration}{The adjustment of a model's raw output so that it can
  be read as a probability. Distinct from accuracy: a model can be accurate and still be
  systematically overconfident.}

\newterm{forgetting}{Catastrophic forgetting}{The loss of previously learned performance
  when a trained network is adapted to new data. Observed here when a cardiac model
  fine-tuned on an external dataset lost accuracy on its original one.}

\newterm{leakage}{Data leakage}{Any route by which information from the evaluation set
  reaches the training process. It inflates measured performance without improving the
  model, and is therefore invisible in the metrics meant to detect failure.}

\newterm{dice}{Dice coefficient}{A measure of overlap between a predicted segmentation and
  its reference, ranging from 0 (disjoint) to 1 (identical).}

\newterm{digitaltwin}{Digital twin}{A simulated model of a patient allowing candidate
  decisions to be explored before one is committed to. Outside the scope of this
  internship and discussed only as a perspective.}

\newterm{escalation}{Escalation}{The decision to withhold a direct answer and route a case
  for further examination, taken when evidence is insufficient, confidence is low, or
  sources disagree.}

\newterm{groupedcv}{Grouped cross-validation}{A partitioning scheme in which every image
  from the same patient or case is placed in the same fold, so that no fold is evaluated
  on a patient it was trained on.}

\newterm{statunit}{Statistical unit}{The entity that may be treated as an independent
  observation. In these datasets it is the patient or the case, not the individual frame,
  because frames from one acquisition are near-duplicates of one another.}

\newterm{triage}{Triage}{The initial assessment that assigns an arriving patient a
  priority level, determining how quickly they are seen.}
